\bookchapter{פרק 22: כיצד AI עובד בטלפון}

\chapterintro{צילום חישובי AI של היום דומה יותר לדחיפת צוות מקצועי של רטושר וצלם ישירות לתוך כפתור הצילום. לפני שאתה לוחץ על הצילום, AI כבר מתחיל לעבוד: סורק את הסצנה, מזהה “זה שמיים כחולים”, “זה פנים אנושיות”, “זה צילום לילה”.}

\booksection{אחד, מצלמת AI: “חדר החושך לאחר עיבוד” עבר לטלפון}

שומע בבירור (זיהוי קולי)

\booksection{שניים, עוזר קולי: כיצד “המזכיר הקולי” שלך מבין אותך?}

\begin{enumerate}
  \item \textbf{מבין (הבנת שפה טבעית)}: מבין כוונה, מחלץ ישויות מפתח.
  \item \textbf{מבצע ועונה}: קורא ל-App או שירות מתאים, ואז "אומר" לך דרך סינתזה קולית.
  \item \textbf{שלוש, AI במכשיר לעומת AI בענן}AI במכשיר משאיר את החישוב בטלפון, מהיר, פרטי יותר, עובד ללא חיבור. AI בענן דורש חיבור, אך יכול לטפל במשימות מורכבות יותר.
\end{enumerate}

\booksection{שלוש, AI במכשיר לעומת AI בענן}

ב-2024-2025, התפוצצות מודלי AI גדולים (כמו DeepSeek, ChatGPT, Wenxin Yiyan, Tongyi Qianwen וכו') שינתה בעומק את אופן השימוש בטלפון. הטלפון כבר לא רק כלי לביצוע הוראות, אלא שותף חכם שיכול להבין, לחשוב וליצור.

\booksection{ארבע, מודלי AI גדולים וטלפון: העוזר החכם האישי שלך}

2025-2026, AI (DeepSeek R1, ChatGPT-5, Wenxin 5.0, Qwen 3) שינה את הטלפון. לא רק כלי, אלא פרטנר. NPU 80+ TOPS = מודל מקומי. “AI Phone” ממחיר שיווקי למוצר אמיתי. NPU מאפשר הרצה מקומית בלי רשת.

\booksubsection{יישומי מודלי AI גדולים בטלפון:}

עוזר שיחה

\booksubsection{: אפליקציות כמו DeepSeek, ChatGPT, Doubao, Kimi, יכולות לענות על שאלות, לתרגם, לכתוב טקסטים, לסכם מאמרים בכל עת. חכם יותר ממנועי חיפוש מסורתיים — הוא לא רק נותן לך קישורים, אלא ישירות תשובות.}

\begin{itemize}
  \item \textbf{כתיבה ויצירה}: תן ל-AI לכתוב עבורך מיילים, לסכם פרוטוקולי פגישות, ללטש מאמרים, ליצור טקסטים יצירתיים.
  \item \textbf{יצירה ועריכת תמונות}: AI מסיר בלחיצה אנשים מתמונות, מרחיב גבולות תמונה, ממיר סגנון, או אפילו יוצר תמונות לפי תיאור טקסט.
  \item \textbf{תרגום בזמן אמת}: תרגום שיחות בזמן אמת, תרגום שיחות פנים אל פנים, תרגום צילום — אין יותר מחסום שפה בנסיעות לחו"ל.
  \item \textbf{עזרה בלימודים}: AI יכול להסביר מושגים, לתרגל, לבדוק שיעורי בית, מורה פרטי נייד.
  \item \textbf{מודל גדול במכשיר: הטלפון מריץ AI בעצמו}טלפוני הדגל האחרונים כבר יכולים להריץ מודלי AI קטנים באופן מקומי (בדרך כלל כ-7B פרמטרים), מה שמביא לשינויים:
\end{itemize}

\booksubsection{לא דורש רשת}

: אפשר להשתמש בפונקציות AI גם ללא חיבור, בטיסה, ברכבת התחתית.

\begin{itemize}
  \item \textbf{פרטיות בטוחה יותר}: השיחות והנתונים שלך לא מועלים לענן, מעובדים כולם מקומית בטלפון.
  \item \textbf{תגובה מהירה יותר}: ללא השהיית רשת, תגובת AI כמעט מיידית.
  \item \textbf{מגבלה נוכחית}: יכולת מודל במכשיר עדיין נחותה ממודל ענן גדול, משימות מורכבות עדיין דורשות חיבור. אבל עם שיפור כוח המחשוב של AI בשבבים, מודלים במכשיר ילכו ויתחזקו.
  \item \textbf{שילוב AI ברמת מערכת:}Apple Intelligence של אפל
\end{itemize}

\booksubsection{: מערכת AI שהוצגה ב-iOS 18, משלבת כלי כתיבה, שרביט קסם לתמונות, Siri משופר, סיכום התראות ועוד. עיבוד במכשיר תחילה, משימות מורכבות עולות לענן.}

\begin{itemize}
  \item \textbf{מודל Pangu הגדול של Huawei}: משולב במערכת HarmonyOS, העוזר Xiaoyi מקבל יכולות מודל גדול, תומך בסיכום מסמכים, זיהוי מסך חכם, עריכת תמונות AI ועוד.
  \item \textbf{Pengpai AI של Xiaomi}: יכולות AI ברמת מערכת, תומך בחיפוש תמונות AI, סיכום AI, כתוביות בזמן אמת AI, שיחה עם מודל גדול במכשיר ועוד.
  \item \textbf{Galaxy AI של Samsung}: תרגום שיחות בזמן אמת, סיכום פתקים AI, יצירה ועריכת תמונות, חיפוש Circle to Search ועוד.
  \item \textbf{חמישה, כמה תרחישי שימוש אופייניים ותרגול}כדי לתת לך להרגיש באופן אינטואיטיבי את "הכוח-על" של AI בטלפון, להלן חמישה תרחישים ספציפיים, המציגים כיצד AI "נחלץ לעזרה" בחיי היומיום:
\end{itemize}

\booksection{כמה תרחישי שימוש טיפוסיים ותרגול מעשי}

ביצועים בפועל

\booksubsection{: מזמין מנה בפריז ולא יודע צרפתית? אפליקציית “תרגום” של iPhone תומכת במצב שיחה דו-כיווני פנים אל פנים, שני הצדדים רואים כתוביות בשפה שלהם, ואפילו אפשר לתרגם בו-זמנית. במחנה Android, תרגום שיחות Galaxy AI של Samsung מאפשר לך לדבר סינית בטלפון והצד השני שומע אנגלית, ולהפך.}

\textbf{העיקרון מאחורה}: זיהוי קולי (ASR) תרגום מכונה (NMT) סינתזה קולית (TTS) פס ייצור בן שלושה שלבים, השהייה מקצה לקצה כבר יכולה להיות מתחת לשנייה.

\textbf{שישה, אופטימיזציית פרומפט: “נוסחת התקשורת” שגורמת ל-AI להקשיב לך}ביצועים בפועל

\booksubsection{: צילמת סלפי מהמם באתר תיירות, אבל הרקע צפוף בתיירים אחרים? השתמש ב“הסרה קסומה” וסמן את עוברי האורח, הטלפון יזהה אוטומטית ו“ימחק” אותם, ואז ימלא את הרקע בצורה חכמה, התוצאה לרוב מדהימה.}

\textbf{העיקרון מאחורה}: מודל פילוח תמונה (Inpainting) מזהה ומוחק תחילה את האזור שנבחר, ואז לפי פיקסלים סמוכים והבנה סמנטית “ממציא” תוכן מילוי סביר. מודל ה-AI מאחורי זה אומן על מיליוני תמונות.

\textbf{שישה, אופטימיזציית פרומפט: “נוסחת התקשורת” שגורמת ל-AI להקשיב לך}ביצועים בפועל

\booksubsection{: צילמת תמונת נוף אפורה ביום מעונן? “שמיים קסומים” של AI יכול בלחיצה להחליף את השמיים לשמיים בהירים, שקיעה או אפילו שמיים זרועי כוכבים. הרחבת תמונה AI יכולה להפוך קומפוזיציה אנכית לאופקית, הטלפון “ממציא” את הסצנה מחוץ למסגרת.}

\textbf{העיקרון מאחורה}: מבוסס על מודל דיפוזיה (Diffusion Model) של AI גנרטיבי (Generative AI), לאחר הבנת התוכן הסמנטי של התמונה, מייצר פיקסלים חדשים. זהו “חזית הטכנולוגיה” הנוכחית של עיבוד תמונות AI.

\textbf{שישה, אופטימיזציית פרומפט: “נוסחת התקשורת” שגורמת ל-AI להקשיב לך}ביצועים בפועל

\booksubsection{: לאחר סיום פגישה, זרוק ל-AI את הקלטת הפגישה, והוא יכול אוטומטית לארגן פרוטוקול פגישה מובנה, לחלץ משימות, ליצור סיכום. דוח PDF באנגלית בן 10 עמודים, AI יכול לסכם את המסקנות המרכזיות תוך 10 שניות.}

\textbf{העיקרון מאחורה}: מודל LLM עם הקשר ארוך (Long Context LLM) יכול “לקרוא” מסמכים בני עשרות אלפי מילים, לבצע חשיבה מרובת שלבים וחילוץ מידע. Apple Intelligence, Xiaoyi של Huawei, WPS AI, Tencent Yuanbao ואחרים מסוגלים לכך.

\textbf{שישה, אופטימיזציית פרומפט: “נוסחת התקשורת” שגורמת ל-AI להקשיב לך}ביצועים בפועל

\booksubsection{: רואה ברחוב פרח שלא מזוהה, פותח מצלמה, ו-AI מיד אומר “זהו פרח היביסקוס”. מחפש באלבום תמונות של הים מהשנה שעברה, חפש “ים”, ו-AI מוצא אוטומטית את כל התמונות עם ים, בלי צורך לדפדף ידנית.}

\textbf{העיקרון מאחורה}: מודל ראייה ממוחשבת (CV) מבצע הטמעה (Embedding) של תוכן התמונה, ואז מתאים דמיון עם שאילתת טקסט. “חיפוש רב-מודאלי” זה הוא קפיצת מדרגה בניהול אלבום הטלפון.

\textbf{שישה, אופטימיזציית פרומפט: “נוסחת התקשורת” שגורמת ל-AI להקשיב לך}איכות התשובות של AI תלויה במידה רבה באופן שבו אתה “שואל”. אותה שאלה, בנוסח שונה, תניב תשובות שונות לחלוטין. להלן נציג גישת תקשורת

\booksection{אופטימיזציה של פרומפטים: “נוסחת התקשורת” שתגרום ל-AI להקשיב לכם}

, עם הרבה דוגמאות.\textbf{למה צריך מסגרת?}הרבה אנשים משתמשים ב-AI כמו ב-Baidu — מזינים כמה מילות מפתח. אבל AI אינו מנוע חיפוש, הוא מודל שפה ש“קורא ומבין ואז עונה”.

\begin{tcbinfo}{תן לו יותר הקשר, הסבר מה אתה רוצה, הגבל אותו שלא יסטה}
, והוא ייתן לך תוצאות מדויקות. המסגרת היא לתרגם את “איך נראית תשובה טובה” בראש שלך להוראות ש-AI יכול להבין.\textbf{1. תפקיד (Role): תגיד ל-AI מי הוא}תן ל-AI “זהות” ספציפית, והוא יפעיל את הידע וסגנון הביטוי המקצועי בתחום המתאים.

\end{tcbinfo}

\booksubsection{ככל שהגדרת התפקיד ספציפית יותר, הפלט מקצועי יותר}

\textbf{דוגמאות מעשיות:}

\booksubsection{לא יותר מ-200 מילה, רשום בנקודות}

\begin{itemize}
  \item ." (שאלת בריאות)\textbf{"אתה}מורה לפיזיקה בחטיבת ביניים שמומחה בהסברת מושגים מורכבים בפשטות
  \item ." (שאלת בריאות)\textbf{"אתה}עורך ספרים עם טעם אסתטי טוב, קורא הרבה, בעל סגנון משלו
  \item ." (שאלת בריאות)\textbf{"אתה}בלוגר בישול שמכין ארוחת בוקר לאנשי משרד כל יום
  \item ." (שאלת בריאות)\textbf{"אתה}רופא סיני מסורתי סבלני, עדין, שלעולם לא לועג למתחילים
  \item ." (שאלת בריאות)\textbf{למה זה עובד?}התפקיד מגדיר את
\end{itemize}

\textbf{5. הערכת איכות פרומפט: 5 מאפיינים של פרומפט טוב},\textbf{סגנון הביטוי}ו\textbf{עומק המקצועיות}. אותה שאלה, מוסברת אחרת לתלמיד ולהורה; המלצות שונות לתושב עיר ולתושב עיירה. בלי רקע, AI יכול לתת רק תשובות “המכנה המשותף הגדול ביותר”, לרוב לא משמעותיות.\textbf{2. רקע (Background): תגיד ל-AI מה המצב שלך}הסבר מראש את “מי אתה, מה המצב הנוכחי, מה המגבלות שלך”, ו-AI ייתן תשובות

\booksubsection{שבאמת מתאימות לך}

, במקום הכללות.\textbf{דוגמאות מעשיות:}“הטלפון שלי הוא

\booksubsection{לא יותר מ-200 מילה, רשום בנקודות}

\begin{itemize}
  \item ." (מצב מכשיר)\textbf{"אני}גמלאית בת 60, עם ראייה לא טובה, לא ממש יודעת להשתמש בסמארטפון
  \item ." (פרופיל משתמש)\textbf{"אני}מתכונן לבחינת אנגלית CET-6 ביוני 2026, רמה בסיסית בינונית
  \item ." (רקע לימודי)\textbf{"אני מנהל}בית קפה קטן בעיר מדרגה שלישית, מחזור חודשי 30 אלף, שני עובדים
  \item ." (רקע עסקי)\textbf{"אני}מיד אחרי לידה בחופשת לידה, שינה לקויה ברצינות
  \item ." (שלב חיים)\textbf{למה זה עובד?}מידע הרקע נותן ל-AI לדעת
\end{itemize}

\textbf{5. הערכת איכות פרומפט: 5 מאפיינים של פרומפט טוב}ו\textbf{“באילו תנאים”}. אותה שאלה, מוסברת אחרת לתלמיד ולהורה; המלצות שונות לתושב עיר ולתושב עיירה. בלי רקע, AI יכול לתת רק תשובות “המכנה המשותף הגדול ביותר”, לרוב לא משמעותיות.\textbf{3. משימה (Task): הסבר ל-AI בבירור מה לעשות}בשלב זה השתמש ב

\booksubsection{פועל בתחילת}

(כתוב, רשום, השווה, תרגם, הסבר, סכם...), תאר את המשימה ב\textbf{אופן ספציפי ובר ביצוע}. משימה מעורפלת תניב רק תשובות מעורפלות.\textbf{דוגמאות מעשיות:}“אנא

\booksubsection{לא יותר מ-200 מילה, רשום בנקודות}

\begin{itemize}
  \item , מופנים לאנשי משרד, זוגות, הורים עם ילדים." (חלוקה ברורה)\textbf{"אנא}תרגם את המייל האנגלי הבא לסינית
  \item , מופנים לאנשי משרד, זוגות, הורים עם ילדים." (חלוקה ברורה)\textbf{"אנא}השווה בין ארבעת הטלפונים בטווח 3000 יואן האלה
  \item , מופנים לאנשי משרד, זוגות, הורים עם ילדים." (חלוקה ברורה)\textbf{"אנא}הסבר 'מה זה 5G' במילים שתלמיד בית ספר יבין
  \item , מופנים לאנשי משרד, זוגות, הורים עם ילדים." (חלוקה ברורה)\textbf{"אנא}כתוב שלושה טקסטים לסרטונים קצרים לחנות ארוחת הבוקר שלי
  \item , מופנים לאנשי משרד, זוגות, הורים עם ילדים." (חלוקה ברורה)\textbf{למה זה עובד?}חלק המשימה הוא "מרכז ההוראות" לפעולת AI.
\end{itemize}

\textbf{5. הערכת איכות פרומפט: 5 מאפיינים של פרומפט טוב}. “תסדר”, “תעשה”, “תראה” מעורפלים הם טאבו גדול ל-AI — הוא ינחש את כוונתך, והתוצאה לרוב תסטה.\textbf{4. מגבלה (Constraints): תגיד ל-AI מה לא לעשות}“מגבלה” נשמעת כמו “אילוץ”, אבל למעשה היא

\booksubsection{כלי רב עוצמה לכווץ את הפלט של AI למה שהכי מתאים לך}

. ממדים נפוצים: מספר מילים, פורמט, טון, דברים להימנע מהם, דברים שחובה לכלול.\textbf{דוגמאות מעשיות:}“

\booksubsection{לא יותר מ-200 מילה, רשום בנקודות}

\begin{itemize}
  \item , רק הסבר את העיקרון." (גבול)\textbf{"}הטון יהיה חם וטבעי, כמו שיחה בין חברים
  \item , רק הסבר את העיקרון." (גבול)\textbf{"}אל תשתמש במונחים מקצועיים
  \item , רק הסבר את העיקרון." (גבול)\textbf{"}לא יותר מ-3 שורות בכל פסקה
  \item , רק הסבר את העיקרון." (גבול)\textbf{"}חובה לכלול: תרחיש שימוש, יתרונות וחסרונות, התאמה לתקציב שלי
  \item , רק הסבר את העיקרון." (גבול)\textbf{"}אל תיתן המלצות השקעה ספציפיות
  \item , רק הסבר את העיקרון." (גבול)\textbf{למה זה עובד?}בלי מגבלות, AI נוטה "להשתולל" — או אקדמי מדי, או ארוך מדי, או סוטה מהנושא. מגבלות ברורות = "סטנדרט אספקה" ברור, שני הצדדים מיישרים ציפיות, והתוצאה נשלטת.
\end{itemize}

\textbf{5. הערכת איכות פרומפט: 5 מאפיינים של פרומפט טוב}כתיבת פרומפט טוב אינה עניין חד-פעמי,

\booksubsection{לפני שליחה, בדוק חמישה קריטריונים אלה}

:\textbf{ספציפי ולא מופשט}: לא “כתוב מאמר”, אלא “כתוב מדריך בחירת טלפון בן 800 מילה, 3 סעיפים, לאמהות חדשות”.

\begin{enumerate}
  \item \textbf{כולל רקע ולא מבודד}: תגיד ל-AI "מי אני, למה אני שואל", ורק אז הוא יכול לתת תשובה "מותאמת אישית".
  \item \textbf{פורמט פלט ברור}: טבלה, רשימה, פסקה, או JSON? תגיד מראש.
  \item \textbf{מגבלות במקום}: מילים, סגנון, דברים להימנע מהם, אחד מהם כבר משפר משמעותית.
  \item \textbf{ניתן לאימות על ידי AI}: לאחר השאלה, בקש מ-AI לחזור על הדרישות שלך; אם לא מצליח להסביר, אז לא הסברת טוב.
  \item \textbf{6. בעיות נפוצות ופתרונות}תסמין
\end{enumerate}

\booksubsection{פתרון}

\begin{center}
\begin{tabularx}{\linewidth}{|X|X|}
\hline
\textbf{תשובה כללית מדי, לא שימושית} & \textbf{הוסף רקע (מי אתה/איזה תרחיש)} \tabularnewline
\hline
תשובה אקדמית/רשמית מדי & הוסף מגבלת טון + רשימת סירוב ל"פקידות" \tabularnewline
תשובה סוטה/לא עונה על השאלה & נסח מחדש את המשימה בפועל, ודרוש "ענה רק על זה" \tabularnewline
תשובה ארוכה מדי, לא מסיימים לקרוא & מגבלה קשיחה של מספר מילים (עד 200 מילה) \tabularnewline
AI ממציא ברצינות & הוסף "אם לא בטוח, סמן 'אני לא בטוח'" + בדיקה עצמית של עובדות חשובות \tabularnewline
לא מרוצה מהגרסה הראשונה & אל תכתוב מחדש, הוסף לפרומפט המקורי "עוד קצת..." \tabularnewline
7. שלושה תרחישים מלאים לדוגמה & תרחיש 1: עוזר חיים (המלצות לארוחת בוקר) \tabularnewline
\hline
\end{tabularx}
\end{center}

\booksubsection{שלושה תרחישים מלאים}

\begin{tcbinfo}{תרחיש 1: עוזר לחיי היומיום (המלצות לארוחת בוקר)}
\textbf{פרומפט}:

אתה בלוגר ארוחות בוקר המתמחה בתזונה מאוזנת (\textbf{תפקיד}). אני אמא עובדת, יש לי רק 15 דקות כל בוקר להכין ארוחת בוקר, ובבית ילד בן 6 וחמות בת 60 (\textbf{רקע}). אנא\textbf{רשום שילובי ארוחת בוקר ל-7 ימים ללא חזרות}(\textbf{משימה}). דרישות לכל ארוחת בוקר: (1) זמן הכנה לא יותר מ-15 דקות (2) תזונה מאוזנת (חלבון+פחמימות+ירקות ופירות) (3) מתאים לכל הגילאים (4) מרכיבים זמינים בכל סופר רגיל (\textbf{מגבלה}).

\textbf{למה זה עובד}: התפקיד קובע את רמת המקצועיות, הרקע נותן את המכנה המשותף לשלוש האוכלוסיות, המשימה מגדירה את התוצר, וארבע המגבלות מבטיחות שהתוכנית ניתנת לביצוע.

\end{tcbinfo}

\begin{tcbinfo}{תרחיש 2: סיוע לימודי (הכנה לבחינת CET-6 באנגלית)}
\textbf{פרומפט}:

אתה\textbf{מורה אנגלית מנוסה לבחינת CET-6}(תפקיד). אני סטודנט לתואר ראשון,\textbf{480 ב-CET-4, יעד 500 ב-CET-6, נותרו 3 חודשים}(רקע). אנא\textbf{הכן עבורי תוכנית הכנה ל-12 שבועות שניתן לבצע}(משימה). דרישות: (1) חלוקת זמן שבועית ברורה להאזנה, דיבור, קריאה וכתיבה (2) כיוון פריצת הדרך לכל חודש (3) המלצה על 3 חומרי עזר מתאימים (עם נימוק, בסינית) (4) פלט בפורמט טבלה (מגבלה).\textbf{למה זה עובד}: התפקיד קובע רקע מקצועי, הרקע נותן נקודות עוגן של ציונים וזמן, המשימה מבהירה “תוכנית מובנית”, המגבלה דורשת “טבלה” ו“ניתנת לביצוע”.

\textbf{שימוש מתקדם: הפוך את AI ל“עוזר אישי” שלך}פרומפט

\end{tcbinfo}

\begin{tcbinfo}{תרחיש 3: יעילות בעבודה (סיכום פרוטוקול ישיבה)}
\textbf{פרומפט}:

אתה\textbf{עוזר/ת אדמיניסטרטיבי/ת ותיק/ה}(תפקיד). זה עתה הסתיימה ישיבת מחלקה של שעה,\textbf{בנושא תוכנית השיווק לרבעון Q2}(רקע).\textbf{אנא ארגן/י את רישום הישיבה לפרוטוקול מובנה}(משימה). דרישות: (1) חלוקה ל“החלטות / משימות / סיכונים והתנגדויות” (2) סמן אחראי ומועד יעד לכל משימה (אם לא צוין — “טעון אישור”) (3) עד 600 מילים בסך הכול (4) פלט בפורמט markdown (מגבלה).\textbf{למה זה עובד}: התפקיד קובע נקודת מבט אדמיניסטרטיבית, הרקע נותן את נושא הפגישה, המשימה מבהירה פלט מובנה, 4 המגבלות הופכות את התוצאה לניתנת לשימוש חוזר.

\end{tcbinfo}

\textbf{שימוש מתקדם: הפוך את AI ל“עוזר אישי” שלך}\par

נסה את השימושים המתקדמים האלה, וחווה את הכוח האמיתי של AI:

\begin{itemize}
  \item \textbf{אופטימיזציה בשיחה מרובת סבבים}: כשאינך מרוצה מתשובת AI, אפשר לומר "זה מקצועי מדי, אנא הסבר שוב במילים שתלמיד כיתה יבין".
  \item \textbf{חשיבה שרשרתית (CoT)}: בקש מ-AI "לחשוב צעד אחר צעד", איכות התשובות לרוב גבוהה יותר: "אנא נתח את בעיית המתמטיקה הזו צעד אחר צעד".
  \item \textbf{שיתוף פעולה בין מספר AI}: השתמש ב-A כדי לכתוב טיוטה, B לליטוש, C לתרגום — שילוב נותן תוצאות טובות יותר.
  \item \textbf{פרומפט מערכת}: בהגדרות AI קבע מראש את הזהות, ההעדפות, הטאבו וכו' שלך, כדי שהתשובות יתאימו יותר לצרכים שלך.
  \item \textbf{סיכון לפרטיות}: בשירותי AI בענן, השיחות שלך עשויות לשמש לאימון המודל. למידע רגיש עדיף AI מקומי (על המכשיר).
\end{itemize}

\booksection{מגבלותיו ואתגרי ה-AI: שמרו על שיקול דעת צלול}

AI חזק, אבל בהחלט לא “יכול הכל”. הכרה צלולה של מגבלותיו מאפשרת להשתמש בו טוב יותר:

\begin{itemize}
  \item \textbf{בעיית "הזיות"}: AI עלול לדבר שטויות ברצינות רבה (במונח המקצועי: Hallucination). הוא עלול להמציא אירועים היסטוריים שלא היו, ציטוטים שגויים ונתונים סטטיסטיים מזויפים. מידע חשוב חייב אימות אנושי.
  \item \textbf{קיצוץ ידע}: לרוב מודלי AI יש תאריך קיצוץ לנתוני האימון (כמו תחילת 2024), הם לא יודעים מה קרה אחר כך.
  \item \textbf{יכולת חשיבה מוגבלת}: חשיבה מרובת שלבים מורכבת, הוכחות מתמטיות, מלכודות לוגיות — AI עדיין נוטה לטעויות.
  \item \textbf{סיכון פרטיות}: בשירותי AI בענן, השיחות שלך עלולות לשמש לאימון מודלים. מידע רגיש עדיף לעבד ב-AI במכשיר.
  \item \textbf{הטיות וערכים}: AI עלול לשאת הטיות מנתוני האימון, יש לשמור על חשיבה עצמאית.
  \item \textbf{טיפים לשימוש ב-AI}תוכן ש-AI מייצר עלול להיות לא מדויק (בעיית "הזיות"), מידע חשוב חובה לאמת.
\end{itemize}

\begin{tcbinfo}{אל תזין ל-AI מידע אישי רגיש (מספר תעודת זהות, סיסמאות, פרטי בנק וכו').}
\begin{itemize}
  \item השתמש נכון בפרומפט: ככל שהתיאור ספציפי וברור יותר, איכות התשובה של AI גבוהה יותר.
  \item AI הוא כלי ולא סמכות, שמור על חשיבה עצמאית ושיקול דעת.
  \item בהחלטות קריטיות (רפואיות, משפטיות, פיננסיות), התייעץ עם איש מקצוע ולא עם AI.
  \item שימוש מתקדם: הפוך את AI ל"עוזר אישי" שלך
  \item נסה את השימושים המתקדמים האלה, וחווה את הכוח האמיתי של AI:
\end{itemize}

\end{tcbinfo}

\begin{tcbpractice}{: אמור לעוזר הקולי: "מחר בבוקר ב-7 תזכיר לי לקחת מטריה, ותעביר את פגישת היומן של היום ב-15:00 ל-16:00." ראה אם הוא יכול לבצע פעולות מרובות שלבים בבת אחת.}
\begin{enumerate}
  \item \textbf{אתגר 2: עריכת תמונות AI}: צלם תמונת תיירות עם עוברי אורח, השתמש ב"הסרה קסומה" או "מחק AI", נקה את עוברי האורח. צפה ב"המצאה" של AI לרקע.
  \item \textbf{אתגר 3: אופטימיזציית פרומפט}: תחילה השתמש בפרומפט פשוט "כתוב מדריך טיול" ל-AI, ואז כתוב פרומפט מפורט במסגרת CRISPE. השווה את ההבדל באיכות בין שתי התשובות.
  \item \textbf{פרק 23: חשיבה תכנותית ואוטומציה למתחילים}בפרקים הקודמים בנינו לטלפון "חומות הגנה" מאובטחות, וגם התוודענו ל"מוח ה-AI" החכם שלו שבתוך החומות. עכשיו, בואו נגלה את אחד הפוטנציאלים האולטימטיביים של הטלפון: לאפשר לו "לחשוב" כמו מתכנת, ולבצע באופן אוטומטי משימות רציפות שאתה מגדיר. אולי תאמר: "אני לא יודע לתכנת, מה זה קשור אליי?" אל תדאג, הפרק הזה לא דורש ממך לכתוב שורה אחת של קוד – טלפונים מודרניים כבר הורידו את רף התכנות לרמה שניתן להסתדר עם גרירה ושחרור.
\end{enumerate}

\end{tcbpractice}

\begin{tcbinfo}{מקרה 2: תהליך מלא להעברת נתונים בין מכשירים}
\begin{itemize}
  \item \textbf{תרחיש}:iPhone 13 > Xiaomi 15 חוצה מותג.
  \item \textbf{צעדים}:1) iCloud/WeChat גיבוי; 2) ניתוק Apple ID + כיבוי "מצא"; 3) Xiaomi App; 4) בחר; 5) WeChat; 6) אימות (אנשי קשר/אלבום/WeChat); 7) איפוס מוצפן.
  \item \textbf{תוצאה צפויה}: אפס אובדן נתונים, המכשיר הישן ניתן למכירה בטוחה.
\end{itemize}

\end{tcbinfo}
